% Options for packages loaded elsewhere
\PassOptionsToPackage{unicode}{hyperref}
\PassOptionsToPackage{hyphens}{url}
\documentclass[
]{article}
\usepackage{xcolor}
\usepackage[margin=1in]{geometry}
\usepackage{amsmath,amssymb}
\setcounter{secnumdepth}{-\maxdimen} % remove section numbering
\usepackage{iftex}
\ifPDFTeX
  \usepackage[T1]{fontenc}
  \usepackage[utf8]{inputenc}
  \usepackage{textcomp} % provide euro and other symbols
\else % if luatex or xetex
  \usepackage{unicode-math} % this also loads fontspec
  \defaultfontfeatures{Scale=MatchLowercase}
  \defaultfontfeatures[\rmfamily]{Ligatures=TeX,Scale=1}
\fi
\usepackage{lmodern}
\ifPDFTeX\else
  % xetex/luatex font selection
\fi
% Use upquote if available, for straight quotes in verbatim environments
\IfFileExists{upquote.sty}{\usepackage{upquote}}{}
\IfFileExists{microtype.sty}{% use microtype if available
  \usepackage[]{microtype}
  \UseMicrotypeSet[protrusion]{basicmath} % disable protrusion for tt fonts
}{}
\makeatletter
\@ifundefined{KOMAClassName}{% if non-KOMA class
  \IfFileExists{parskip.sty}{%
    \usepackage{parskip}
  }{% else
    \setlength{\parindent}{0pt}
    \setlength{\parskip}{6pt plus 2pt minus 1pt}}
}{% if KOMA class
  \KOMAoptions{parskip=half}}
\makeatother
\usepackage{color}
\usepackage{fancyvrb}
\newcommand{\VerbBar}{|}
\newcommand{\VERB}{\Verb[commandchars=\\\{\}]}
\DefineVerbatimEnvironment{Highlighting}{Verbatim}{commandchars=\\\{\}}
% Add ',fontsize=\small' for more characters per line
\usepackage{framed}
\definecolor{shadecolor}{RGB}{248,248,248}
\newenvironment{Shaded}{\begin{snugshade}}{\end{snugshade}}
\newcommand{\AlertTok}[1]{\textcolor[rgb]{0.94,0.16,0.16}{#1}}
\newcommand{\AnnotationTok}[1]{\textcolor[rgb]{0.56,0.35,0.01}{\textbf{\textit{#1}}}}
\newcommand{\AttributeTok}[1]{\textcolor[rgb]{0.13,0.29,0.53}{#1}}
\newcommand{\BaseNTok}[1]{\textcolor[rgb]{0.00,0.00,0.81}{#1}}
\newcommand{\BuiltInTok}[1]{#1}
\newcommand{\CharTok}[1]{\textcolor[rgb]{0.31,0.60,0.02}{#1}}
\newcommand{\CommentTok}[1]{\textcolor[rgb]{0.56,0.35,0.01}{\textit{#1}}}
\newcommand{\CommentVarTok}[1]{\textcolor[rgb]{0.56,0.35,0.01}{\textbf{\textit{#1}}}}
\newcommand{\ConstantTok}[1]{\textcolor[rgb]{0.56,0.35,0.01}{#1}}
\newcommand{\ControlFlowTok}[1]{\textcolor[rgb]{0.13,0.29,0.53}{\textbf{#1}}}
\newcommand{\DataTypeTok}[1]{\textcolor[rgb]{0.13,0.29,0.53}{#1}}
\newcommand{\DecValTok}[1]{\textcolor[rgb]{0.00,0.00,0.81}{#1}}
\newcommand{\DocumentationTok}[1]{\textcolor[rgb]{0.56,0.35,0.01}{\textbf{\textit{#1}}}}
\newcommand{\ErrorTok}[1]{\textcolor[rgb]{0.64,0.00,0.00}{\textbf{#1}}}
\newcommand{\ExtensionTok}[1]{#1}
\newcommand{\FloatTok}[1]{\textcolor[rgb]{0.00,0.00,0.81}{#1}}
\newcommand{\FunctionTok}[1]{\textcolor[rgb]{0.13,0.29,0.53}{\textbf{#1}}}
\newcommand{\ImportTok}[1]{#1}
\newcommand{\InformationTok}[1]{\textcolor[rgb]{0.56,0.35,0.01}{\textbf{\textit{#1}}}}
\newcommand{\KeywordTok}[1]{\textcolor[rgb]{0.13,0.29,0.53}{\textbf{#1}}}
\newcommand{\NormalTok}[1]{#1}
\newcommand{\OperatorTok}[1]{\textcolor[rgb]{0.81,0.36,0.00}{\textbf{#1}}}
\newcommand{\OtherTok}[1]{\textcolor[rgb]{0.56,0.35,0.01}{#1}}
\newcommand{\PreprocessorTok}[1]{\textcolor[rgb]{0.56,0.35,0.01}{\textit{#1}}}
\newcommand{\RegionMarkerTok}[1]{#1}
\newcommand{\SpecialCharTok}[1]{\textcolor[rgb]{0.81,0.36,0.00}{\textbf{#1}}}
\newcommand{\SpecialStringTok}[1]{\textcolor[rgb]{0.31,0.60,0.02}{#1}}
\newcommand{\StringTok}[1]{\textcolor[rgb]{0.31,0.60,0.02}{#1}}
\newcommand{\VariableTok}[1]{\textcolor[rgb]{0.00,0.00,0.00}{#1}}
\newcommand{\VerbatimStringTok}[1]{\textcolor[rgb]{0.31,0.60,0.02}{#1}}
\newcommand{\WarningTok}[1]{\textcolor[rgb]{0.56,0.35,0.01}{\textbf{\textit{#1}}}}
\usepackage{graphicx}
\makeatletter
\newsavebox\pandoc@box
\newcommand*\pandocbounded[1]{% scales image to fit in text height/width
  \sbox\pandoc@box{#1}%
  \Gscale@div\@tempa{\textheight}{\dimexpr\ht\pandoc@box+\dp\pandoc@box\relax}%
  \Gscale@div\@tempb{\linewidth}{\wd\pandoc@box}%
  \ifdim\@tempb\p@<\@tempa\p@\let\@tempa\@tempb\fi% select the smaller of both
  \ifdim\@tempa\p@<\p@\scalebox{\@tempa}{\usebox\pandoc@box}%
  \else\usebox{\pandoc@box}%
  \fi%
}
% Set default figure placement to htbp
\def\fps@figure{htbp}
\makeatother
\setlength{\emergencystretch}{3em} % prevent overfull lines
\providecommand{\tightlist}{%
  \setlength{\itemsep}{0pt}\setlength{\parskip}{0pt}}
\usepackage{bookmark}
\IfFileExists{xurl.sty}{\usepackage{xurl}}{} % add URL line breaks if available
\urlstyle{same}
\hypersetup{
  pdftitle={R-hat and ESS Diagnostics},
  hidelinks,
  pdfcreator={LaTeX via pandoc}}

\title{R-hat and ESS Diagnostics}
\author{}
\date{\vspace{-2.5em}2026-04-17}

\begin{document}
\maketitle

\subsection{Introduction}\label{introduction}

An MCMC posterior summary is only as good as the chain that produced it.
Two questions decide whether the draws are usable: have the chains
converged on the same distribution, and how many independent
observations is the autocorrelated chain actually worth? The first is
answered by \(\hat R\), the second by effective sample size (ESS). These
two numbers, reported alongside every parameter in modern Bayesian
output, are the minimum diagnostic checklist before any posterior
interval is interpreted.

\subsection{Prerequisites}\label{prerequisites}

A working understanding of Markov chain Monte Carlo, the role of warm-up
versus sampling iterations, and the autocorrelation structure of MCMC
draws.

\subsection{Theory}\label{theory}

\(\hat R\) compares the variance between chains to the variance within
chains. If chains have converged on the same target, the two variance
estimates agree and \(\hat R \to 1\); persistent disagreement inflates
\(\hat R\). The Vehtari et al.~(2021) rank-normalised variant handles
heavy-tailed posteriors and is the default in
\texttt{posterior::summarise\_draws()}. ESS expresses the variance of
the Monte Carlo estimator as if the draws were independent:
\(n_{\text{eff}} = N / (1 + 2 \sum_k \rho_k)\) for autocorrelations
\(\rho_k\). Bulk ESS quantifies precision of the posterior mean and
median; tail ESS quantifies precision of the 5 \% and 95 \% quantiles,
which require many more draws because they are estimated from a thin
part of the distribution.

\subsection{Assumptions}\label{assumptions}

At least four chains started from dispersed initial values, with
adequate warm-up so that the sampler has tuned step size and mass
matrix. A single chain cannot diagnose convergence, regardless of how
long it is --- between-chain variance is undefined.

\subsection{R Implementation}\label{r-implementation}

\begin{Shaded}
\begin{Highlighting}[]
\FunctionTok{library}\NormalTok{(brms); }\FunctionTok{library}\NormalTok{(posterior)}

\FunctionTok{data}\NormalTok{(sleepstudy, }\AttributeTok{package =} \StringTok{"lme4"}\NormalTok{)}
\NormalTok{fit }\OtherTok{\textless{}{-}} \FunctionTok{brm}\NormalTok{(Reaction }\SpecialCharTok{\textasciitilde{}}\NormalTok{ Days }\SpecialCharTok{+}\NormalTok{ (Days }\SpecialCharTok{|}\NormalTok{ Subject), }\AttributeTok{data =}\NormalTok{ sleepstudy,}
           \AttributeTok{chains =} \DecValTok{4}\NormalTok{, }\AttributeTok{iter =} \DecValTok{2000}\NormalTok{, }\AttributeTok{refresh =} \DecValTok{0}\NormalTok{)}

\CommentTok{\# Summary with R{-}hat and ESS}
\FunctionTok{summary}\NormalTok{(fit)}

\CommentTok{\# Per{-}parameter}
\NormalTok{post\_draws }\OtherTok{\textless{}{-}} \FunctionTok{as\_draws\_df}\NormalTok{(fit)}
\FunctionTok{summarise\_draws}\NormalTok{(post\_draws)}
\end{Highlighting}
\end{Shaded}

\subsection{Output \& Results}\label{output-results}

\texttt{summary(fit)} reports \texttt{Rhat}, \texttt{Bulk\_ESS}, and
\texttt{Tail\_ESS} for every parameter. \texttt{summarise\_draws()}
returns the same diagnostics in a tidy tibble, ready for \texttt{dplyr}
filtering when a model has hundreds of parameters and you want to know
which one needs attention.

\subsection{Interpretation}\label{interpretation}

A reporting sentence: ``All parameters had \(\hat R \le 1.00\), bulk ESS
\(> 2{,}000\), and tail ESS \(> 1{,}500\); sampling is adequate for
posterior intervals.'' Conventional thresholds: \(\hat R < 1.01\)
indicates convergence, \(\hat R > 1.05\) is a clear failure; bulk and
tail ESS should each exceed 400 for credible interval estimation to be
reliable.

\subsection{Practical Tips}\label{practical-tips}

\begin{itemize}
\tightlist
\item
  If \(\hat R\) exceeds 1.01, double the iteration count first; if the
  problem persists, reparameterise (e.g.~non-centred random effects)
  rather than throwing more iterations at it.
\item
  Low ESS relative to total draws (ratios under 0.1) signals strong
  autocorrelation; raise \texttt{adapt\_delta} or examine the model for
  funnel-shaped geometry.
\item
  Divergent transitions in HMC are a separate, more serious diagnostic
  --- they indicate that the sampler missed regions of the posterior,
  not just that it explored slowly. Always inspect them with
  \texttt{bayesplot::mcmc\_parcoord()} or \texttt{pairs()}.
\item
  Trace plots remain a useful supplement: a ``hairy caterpillar'' with
  chains overlapping is what convergence looks like; a chain wandering
  away or stuck in a valley is visible long before \(\hat R\) catches it
  for short runs.
\item
  For very high-dimensional models, monitor \(\hat R\) and ESS for
  \emph{derived} quantities of interest (e.g.~predictions, contrasts),
  not just raw parameters; convergence of one does not imply the other.
\end{itemize}

\subsection{R Packages Used}\label{r-packages-used}

\texttt{posterior} for the rank-normalised \(\hat R\), bulk and tail
ESS, and tidy-draw structures, \texttt{bayesplot} for diagnostic plots,
and \texttt{brms} (or \texttt{rstanarm}, \texttt{cmdstanr}) as the
modelling backend that produces the draws in the first place.

\subsection{For Reviewers}\label{for-reviewers}

\textbf{What to look for in a paper using this method.}

\begin{itemize}
\tightlist
\item
  \textbf{Common misapplications.}
\item
  \textbf{Diagnostics that should be reported but often aren't.}
\item
  \textbf{Red flags in tables and figures.}
\item
  \textbf{What to verify.}
\item
  \textbf{What an adequate Methods paragraph must contain.}
\end{itemize}

\subsection{See also --- labs in this
chapter}\label{see-also-labs-in-this-chapter}

\begin{itemize}
\tightlist
\item
  \href{labs/lab-bayesian-thinking-control-vs-decision-error.Rmd}{Bayesian
  thinking; α control vs decision error}
\item
  \href{labs/lab-brms-stan-loo-hierarchical-models.Rmd}{brms / Stan,
  LOO, hierarchical models}
\end{itemize}

\end{document}
